\subsection{APPLICATIONS: III. GENERALIZED POWERS}

Let E be a monoid with identity element e and law of composition denoted by T. If x is invertible in E, let \(g_{x}\) be the homomorphism of Z into E mapping 1 to x. Let \(g_{x}(n) = \frac{n}{T} x\) for all \(n \in Z\) ; by Lemma 2 this notation is compatible for \(n \in N\) with the notation introduced earlier. Then

\[
{ } ^ { m + n } \top x = \left( \begin{array} { c } m \\ T \end{array} x \right) \top \left( \begin{array} { c } n \\ T \end{array} x \right)\tag{9}
\]

\[
\frac {0}{T} x = e\tag{10}
\]

\[
\frac {1}{T} x = x\tag{11}
\]

for \(x\) invertible in \(\mathbf{E}\) and \(m, n\) in \(\mathbf{Z}\) . Further, if \(y = \frac{m}{T} x\) , the mapping \(n \mapsto g_x(mn)\) of \(\mathbf{Z}\) into \(\mathbf{E}\) is a homomorphism mapping 1 to \(y\) , whence \(g_x(mn) = g_y(n)\) , that is

\[
\frac {m n}{T} x = \frac {n}{T} \left(\frac {m}{T} x\right).\tag{12}
\]

As \(-1\) is the negative of \(1\) in \(\mathbf{Z}\) , \(\overline{\top}^{-1}x\) is the inverse of \(x = \frac{1}{\top} x\) in E. If we write \(n = -m\) in (9), it is seen that \(\overline{\top}^{-m}x\) is the inverse of \(\overline{\top}^{m}x\) .

